\documentclass[twocolumn]{aastex631}

\usepackage{amsmath}
\usepackage{multirow}
\usepackage{natbib}
\usepackage{graphicx} 
\usepackage{aas_macros}

\begin{document}

Scalar-tensor theories of gravity represent a cornerstone of modern theoretical cosmology, providing a rich framework for addressing phenomena such as cosmic inflation and the accelerated expansion of the universe. Within this landscape, Degenerate Higher-Order Scalar-Tensor (DHOST) theories are particularly noteworthy. These theories are constructed to contain higher-order derivatives of the fields, which would generically introduce pathological ghost instabilities, rendering the theory physically inviable. However, DHOST theories circumvent this fate through a precise tuning of their defining functions, which enforces a degeneracy condition at the level of the Lagrangian. This condition ensures that the would-be ghost degree of freedom is constrained, leaving a classically stable theory with the correct number of propagating modes. This intricate construction, however, raises a foundational question: is this life-saving degeneracy an accidental mathematical feature, or is it the consequence of a deeper, underlying symmetry principle?A compelling hypothesis suggests that the robust structure of DHOST theories is protected by a local, field-dependent symmetry. Unlike standard gauge symmetries, this conjectured transformation nontrivially mixes the scalar and metric sectors. Specifically, for a local spacetime-dependent parameter $\epsilon(x)$, the fields are proposed to transform as\begin{equation}\delta_\epsilon \phi = \epsilon(x) \Lambda(\phi, X) \quad \text{and} \quad \delta_\epsilon g_{\mu\nu} = \epsilon(x) \mathcal{L}_{\xi} g_{\mu\nu}, \citep{popławski2024classicalphysicsspacetimefields,kolev2024objectiveratescovariantderivatives}\end{equation}where $\mathcal{L}_{\xi}$ is the Lie derivative along a vector field $\xi^\mu$ constructed from the scalar field gradient, $\xi^\mu = \alpha(\phi, X)\partial^\mu\phi$, and $X = -\frac{1}{2}g^{\mu\nu}\partial_\mu\phi\partial_\nu\phi$ is the scalar's kinetic term  \citep{gao2019higherderivativescalartensortheory,apostolopoulos2023constructingfamilyconformallyflat}. The functions $\Lambda$ and $\alpha$ depend on the scalar field and its kinetic energy. If this symmetry is indeed the guardian of classical stability, it must be a fundamental feature of the theory, one that should persist beyond the classical approximation.The central challenge, which this paper confronts directly, arises when considering the theory within the framework of effective field theory. Any classical action is understood to be the leading-order approximation to a more complete quantum description. Integrating out high-energy physics is expected to generate an infinite tower of higher-dimension operators, with the leading corrections typically appearing as higher-order curvature invariants. The inclusion of such terms, like the Gauss-Bonnet invariant and the square of the Weyl tensor, provides a crucial test for the proposed symmetry. The core problem is the potential conflict between the highly specific structure required by the symmetry and the generic form of these quantum-motivated corrections. The difficulty of testing this is twofold: the DHOST Lagrangian itself is algebraically complex \citep{langlois2018degeneratehigherorderscalartensordhost}, and the proposed symmetry transformation is non-standard, making a direct verification computationally demanding.In this work, we perform a direct, first-principles investigation into the fate of this conjectured symmetry in the presence of leading-order curvature corrections. We construct an effective action composed of a classically ghost-free DHOST Lagrangian \citep{babichev2023rotatingblackholesembedded,brax2025primordialgravitationalwavesdhost} supplemented by constant-coefficient Gauss-Bonnet and Weyl-squared terms \citep{cisterna2025thermodynamicsfourdimensionalregularblack}. Our method is a rigorous and explicit calculation of the variation of this complete action under the aforementioned symmetry transformation. Crucially, we make no prior assumption of invariance. Instead, we systematically compute the transformation of every constituent term in the Lagrangian to determine precisely how the action responds. For the action to be symmetric, its variation must vanish identically for any arbitrary function $\epsilon(x)$, which, after integration by parts, imposes a stringent differential condition on the functions $\Lambda$ and $\alpha$ that define the transformation.Our analysis provides a definitive resolution to this question. We find a stark dichotomy: the classical DHOST sector of the action is indeed invariant, provided the transformation functions $\Lambda$ and $\alpha$ satisfy a specific differential relation, thus confirming the conjecture at the classical level. However, we demonstrate unequivocally that the standard constant-coefficient Gauss-Bonnet and Weyl-squared terms are not invariant and explicitly break the symmetry. This result establishes a fundamental tension between the symmetry responsible for classical health and the standard form of leading-order quantum corrections. It implies that a self-consistent, quantum-corrected DHOST theory cannot be constructed by naively appending generic curvature invariants; rather, the corrections themselves must be incorporated in a more subtle, symmetry-preserving fashion, providing a critical guideline for the construction of viable quantum theories of modified gravity.

\end{document}
                